\section{Signals and empirical feature selection}
\label{sec:signals}

\subsection{Recorded views}
The feature family contains a normalized within-layer exponential moving average
(EMA) of attention
$x_{\rm wl}=\bar a_{\ell,b,t}/\max_{b'}\bar a_{\ell,b',t}$, a previous-layer
hot-block indicator $x_{\rm cl}=\mathbf 1[\bar a_{\ell-1,b,t}\text{ is in the top-}r]$,
a query-derived proxy $x_{\rm q}=(1+\cos(\tilde q_t,\tilde K_{\ell,b}))/2$ and a
temporal-age feature $x_{\rm age}=\exp(-(t-p_b)/w)$.
Here $\bar a$ is an EMA and $p_b$ the block-creation step, so $x_{\rm age}$ is an
age proxy, not an LRU access-recency measurement; the tildes mark the archived
query representation, whose coordinate-space audit in \S\ref{sec:limits} limits
conclusions about that view. The cross-layer signal is an attention-magnitude
carry-over, not InfiniGen's query-based mechanism, and the offline extractor falls
back to the current layer's own EMA at layer 0 where runtime realizations need
explicit fallback and alignment rules (Table~\ref{tab:protocol}).

\subsection{Relevance diagnostics}
\label{sec:redundancy}
We estimate relevance $I(X_i;Z)$ with quantile-binned plug-in estimators, in nats.
On the version-2 corpus, within-layer EMA and the cross-layer indicator have $I(X_i;Z)$
of 0.104 and 0.099 nats, versus 0.007 for the cosine query proxy and under 0.001 for age
(\S\ref{sec:exp-relevance}); raw-probe AUCs are in Table~\ref{tab:v2}.

Archived conditional-information diagnostics $U_i=I(X_i;Z\mid X_{-i})$ agree that
the magnitude views carry what the proxies do not ($\le0.006$ nats): within-layer
is $0.133$--$0.155$ on every binning grid and the binary cross-layer estimate
$0.048$--$0.084$ on the finer ones, collapsing to zero under the coarsest, so the
artifact reports its grid sensitivity. These estimates motivate a candidate subset,
evaluated on held-out requests; they are not a sufficiency certificate, and the
evidence for the compact model is its measured predictive performance.
